\section{非线性振动与摄动（选读）}\label{ch10:sec:nonlinear}
保留势能的高阶项后，回复力不再线性，叠加原理也不能原样使用。一个保留左右对称性的例子是杜芬方程（Duffing equation）
\[
 \ddot x+\omega_0^2x+\varepsilon\alpha x^3=0,
 \qquad U=\frac12m\omega_0^2x^2+\frac14m\varepsilon\alpha x^4.
\]
$\varepsilon$用来标记小修正阶数，要求$|\varepsilon\alpha A^2|\ll\omega_0^2$。若仍把频率固定为$\omega_0$，摄动方程会产生$t\sin\omega_0t$这样的长期项，时间稍长就失去一致性。为避免它，令$\tau=\Omega t$，$\Omega=\omega_0+\varepsilon\Omega_1+\cdots$，$x=A\cos\tau+\varepsilon x_1(\tau)+\cdots$。

因为$\ddot x=\Omega^2\dd^2x/\dd\tau^2$，收集一阶项得到
\[
 \omega_0^2(x_1''+x_1)=2\omega_0\Omega_1 A\cos\tau-\alpha A^3\cos^3\tau.
\]
利用$\cos^3\tau=(3\cos\tau+\cos3\tau)/4$，右端的基频系数是$2\omega_0\Omega_1A-3\alpha A^3/4$。令它为零消去共振驱动，得到
\[
 \Omega_1=\frac{3\alpha A^2}{8\omega_0},\qquad
 \boxed{\Omega\simeq\omega_0\left(1+\frac{3\varepsilon\alpha A^2}{8\omega_0^2}\right)}.
\]
余下$x_1''+x_1=-\alpha A^3\cos3\tau/(4\omega_0^2)$，特解为$\alpha A^3\cos3\tau/(32\omega_0^2)$。若把$A$定义为$t=0$的转向位移，还应加齐次项使$x_1(0)=0$，从而
\[
 x\simeq A\cos\tau+
 \frac{\varepsilon\alpha A^3}{32\omega_0^2}(\cos3\tau-\cos\tau).
\]
这明确显示三次谐波。$\varepsilon\alpha>0$为硬化，振幅增大时频率升高；$\varepsilon\alpha<0$为软化，只能在平衡附近和模型有效范围讨论，不能外推到负四次势能导致的远处失稳区域。令非线性参数趋于零，以上结果恢复简谐振动。

另一类现象来自参数调制，例如
\[
 \ddot x+\omega_0^2[1+h\cos(\Omega t)]x=0,\qquad |h|\ll1.
\]
外界改变的是回复系数，而不是直接加同频外力。在$\Omega=2\omega_0$附近，设$x=a(t)\cos\omega_0t+b(t)\sin\omega_0t$，$a,b$缓变，忽略其二阶时间导数。保留右端乘积中的基频项，得到
\[
 2\omega_0\dot b=-\frac12h\omega_0^2a,
 \qquad -2\omega_0\dot a=\frac12h\omega_0^2b.
\]
即$\dot a=-h\omega_0b/4$、$\dot b=-h\omega_0a/4$，两种包络组合有指数率$\pm h\omega_0/4$；一般初始微扰包含增长分量。这是主参数共振的近似机制。若另有弱黏性阻尼，增长须克服$\gamma$，在精确调谐的一阶近似下要求$|h|\omega_0/4>\gamma$；频率失谐也会缩小不稳定区。严格零位移、零速度仍是齐次方程的零解，不会在没有扰动时凭空起振。
\subsection{小修正为什么会在长时间后变得重要}
摄动不是把某个小项直接删去，而是在一个明确的小参数下，追踪它对主运动的首阶影响。以本节的三次回复项为例，在振幅 $A$ 附近，非线性加速度与线性加速度的比为 $|\varepsilon\alpha|A^2/\omega_0^2$。只有这个无量纲比值很小，才允许把运动看成简谐振动加小修正。

如果固定原频率写 $x=x_0+\varepsilon x_1$，其中 $x_0=A\cos\omega_0t$，一阶方程会是
\[
 \ddot x_1+\omega_0^2x_1
 =-\frac{3\alpha A^3}{4}\cos\omega_0t
  -\frac{\alpha A^3}{4}\cos3\omega_0t.
\]
第一项恰以系统本征频率驱动修正量。验证
\[
 \left(\frac{\dd^2}{\dd t^2}+\omega_0^2\right)
 [t\sin\omega_0t]=2\omega_0\cos\omega_0t,
\]
便知其中包含 $-3\alpha A^3t\sin\omega_0t/(8\omega_0)$。即使 $\varepsilon$ 很小，乘上足够长的时间后也可能不再小，这就是长期项。
它不表示保守振子真的无界增长，而是说明选错了比较基准。若真实频率为 $\omega_0+\delta\omega$，则
\[
 A\cos[(\omega_0+\delta\omega)t]
 \simeq A\cos\omega_0t-A\delta\omega\,t\sin\omega_0t.
\]
相位误差积累就产生同样形式的长期项。把频率也纳入待求量，是把这部分相位变化吸收到主运动中，而不是人为删除真实物理效应。

\begin{origexamplebox}{弱硬化弹簧的频率与适用范围}\label{ch10:ex:teaching-hardening}
质点质量为 $m$，在光滑直线上受到回复力 $F=-kx-\lambda x^3$，$k,\lambda>0$。从 $x=A$ 静止释放，且 $\lambda A^2/k\ll1$。求振动频率的首阶修正；若要求频率相对简谐值的修正不超过 $1\%$，振幅应满足什么近似条件？
\end{origexamplebox}
\begin{xesolution}
由牛顿定律，$\ddot x+(k/m)x+(\lambda/m)x^3=0$，所以 $\omega_0=\sqrt{k/m}$，而本节的一阶系数对应 $\varepsilon\alpha=\lambda/m$。代入已推导的消长期项条件，得到
\[
 \Omega\simeq\sqrt{\frac km}\left(1+\frac{3\lambda A^2}{8k}\right),
 \qquad T\simeq2\pi\sqrt{\frac mk}\left(1-\frac{3\lambda A^2}{8k}\right).
\]
周期使用 $(1+q)^{-1}\simeq1-q$，不能照抄频率的正号。回复力在较大位移处比线性弹簧更强，频率升高、周期缩短与物理判断相符。
首阶相对修正为 $3\lambda A^2/(8k)$，故近似要求
\[
 A\lesssim\sqrt{\frac{0.08k}{3\lambda}}.
\]
该边界使 $\lambda A^2/k\lesssim0.0267$，与小参数假设一致，但这是首阶误差估计，不是包含全部高阶项的严格误差界。由 $[\lambda]=\mathrm{N/m^3}$ 可知 $k/\lambda$ 为长度平方；$A\to0$ 时恢复线性结果。
\end{xesolution}

\subsection{参数共振的能量从哪里来}
受迫振动是外界施加随时间变化的力；参数共振则是外界改变弹簧劲度、摆长等系统参数。对 $k(t)=m\omega_0^2[1+h\cos(\Omega t)]$，瞬时机械能写作
\[
 E=\frac12m\dot x^2+\frac12k(t)x^2.
\]
对时间求导并代入 $m\ddot x+k(t)x=0$，得到
\[
 \dot E=\dot x[m\ddot x+k(t)x]+\frac12\dot k(t)x^2
 =\frac12\dot k(t)x^2.
\]
所以能量由调节参数的外界装置提供，并未违反守恒。若在位移大时增大劲度、位移小时减小劲度，平均上可以输入能量。一次振动有两次位移绝对值极大，这提示主参数共振出现在约两倍本征频率处，但精确增长率仍须由缓变振幅方程计算。
加入阻力 $-2m\gamma\dot x$ 后，还要减去 $2m\gamma\dot x^2$ 的耗散功率。只有平均输入克服耗散，振幅才增长；超出小振幅范围后，非线性、参数调制范围及装置反馈也会限制增长，线性指数解不能无限外推。

